% Part: turing-machines
% Chapter: machines-computations
% Section: combining-machines

\documentclass[../../../include/open-logic-section]{subfiles}

\begin{document}

\olfileid{tur}{mac}{cmb}
\olsection{د ټيورينګ ماشينونو يوځاے کول}

\begin{explain}
د ټيورينګ ماشينونو هغه بېلګې چې تر اوسه مو ليدلې، نسبتاً ساده وې۔
خو په حقيقت کښې هره مسئله چې په هره اوسنۍ پروګرام‌ليکنې ژبه حل کېدے
شي، د ټيورينګ ماشينونو په وسيله هم حل کېدے شي۔ د لا پېچلو ټيورينګ
ماشينونو د جوړولو دپاره بايد ځان ډاډه کړو چې ماشينونه يوځاے کولے شو؛
په دې توګه به يوه کړنلاره په ساده برخو ووېشو او د لا پېچلو مسئلو د
حل ماشينونه به جوړ کړو۔ که د پېچلې مسئلې په رغوونکو برخو د وېشلو
طبيعي لار پيدا کړو، مسئله په څو پړاوونو کښې حل کولے شو: څو ساده
ټيورينګ ماشينونه جوړوو او بيا ئې په يوه داسې ماشين کښې يوځاے کوو چې
ټوله مسئله حل کړي۔ دا خبره په بله برخه کښې د درېدو مسئلې د څېړلو پر
وخت ځانګړې مهمه ده۔

ټيورينګ ماشينونه $M = \tuple{Q, \Sigma, q_0, \delta}$ او
~$M' = \tuple{Q', \Sigma', q_0', \delta'}$ څنګه يوځاے کړو؟ اوس د
~$M$ له درېدو وروسته د پټې تشکيل د ماشين~$M'$ د يوې چلونې د ننوت
تشکيل په توګه کاروو۔ د داسې يوه ټيورينګ ماشين $M \frown M'$ د لاس
ته راوړلو دپاره لاندې کارونه وکړئ:
\begin{enumerate}
    \item د~$M'$ ټول حالتونه~$Q'$ بيا وشمېرئ (يا بيا ونوموئ)، څو
    $M$ او~$M'$ هېڅ ګډ حالت ونۀ لري ($Q \cap Q' = \emptyset$)۔
    \item د~$M \frown M'$ حالتونه $Q \cup Q'$ دي۔
    \item د پټې الفبا $\Sigma \cup \Sigma'$ ده۔
    \item لومړنی حالت~$q_0$ دے۔
    \item د حالت بدلون تابع هغه تابع~$\delta''$ ده چې داسې ورکول کېږي:
    \[\delta''(q,\sigma) =
    \begin{cases}
      \delta(q,\sigma) & \text{که $q \in Q$ وي او دا ارزښت تعريف شوے وي}\\
      \delta'(q,\sigma) & \text{که $q \in Q'$ وي}\\
      \tuple{q_0', \sigma, \TMstay} & \text{که $q \in Q$ وي او
      $\delta(q,\sigma)$ تعريف شوې نۀ وي}
    \end{cases}\]
\end{enumerate}
لاس ته راغلے ماشين چې په حالت $q \in Q$ کښې وي، د~$M$ لارښوونې
کاروي، او چې په حالت~$q \in Q'$ کښې وي، د~$M'$ لارښوونې کاروي۔ کله
چې دا په داسې حالت $q \in Q$ کښې وي او داسې نښه~$\sigma$ لولي چې
~$M$ ورته د حالت بدلون نۀ لري (يعنې~$M$ به درېدلے و)، د~$M'$ لومړني
حالت ته ننوځي (او د پټې منځپانګه او د سر ځاے هماغسې پرېږدي)۔
\textbf{د ژباړې د نسخې سمون (OLCMP-047):} د سرچينې د ټوټه‌ييز تعريف
لومړي شرط يوازې دا ويل چې حالت د لومړي ماشين دے، او په دې توګه ئې له
درېيم شرط سره ټکر کاوه۔ دلته څرګنده شوه چې لومړی شرط هغه وخت دے چې
د لومړي ماشين ارزښت تعريف شوے وي؛ درېيم شرط ئې ناتعريف شوی حالت
سمبالوي۔

پام وکړئ، تر څو چې ماشين~$M$ منضبط نۀ وي، موږ نۀ پوهېږو چې د~$M$
د درېدو پر وخت د پټې سر چېرته وي؛ نو د~$M$ په درېدني تشکيل کښې سر
اړين نۀ دے چې مربع~$1$ ولولي۔ د ماشينونو د يوځاے کولو پر وخت بايد
دا خبره په پام کښې وساتو۔
\end{explain}

\begin{ex}
\emph{د ماشينونو يوځاے کول:} موږ به داسې ماشين طرح کړو چې د~$n$ او
~$m$ اوږدوالي لرونکو د~$\TMstroke$ نښو له دوو غونډونو جوړ ننوت باندې
پيل شي او پر پټه د~$2(m+n)$ شمېر~$\TMstroke$ نښو په يوه غونډ
ودرېږي۔ د دې ماشين د جوړولو دپاره دوه اشنا ماشينونه يوځاے کولے شو:
د جمع ماشين او دوه‌چنده کوونکے۔ لومړے د جمع ماشين د حالت ډياګرام
کاږو۔
\[
\begin{tikzpicture}[->,>=stealth',shorten >=1pt,auto,node distance=2.8cm,
                    semithick]
  \tikzstyle{every state}=[fill=none,draw=black,text=black]

  \node[initial,state] (A)              {$q_0$};
  \node[state]         (B) [right of=A] {$q_1$};
  \node[state]         (C) [right of=B] {$q_2$};

  \path (A) edge node {\TMtrans{\TMblank}{\TMstroke}{\TMstay}} (B)
            edge [loop above] node {\TMtrans{\TMstroke}{\TMstroke}{\TMright}} (B)
        (B) edge [loop above] node {\TMtrans{\TMstroke}{\TMstroke}{\TMright}} (B)
            edge node {\TMtrans{\TMblank}{\TMblank}{\TMleft}} (C)
        (C) edge [loop above] node {\TMtrans{\TMstroke}{\TMblank}{\TMstay}} (C);
\end{tikzpicture}
\]
موږ نۀ غواړو چې ماشين په حالت~$q_2$ کښې ودرېږي، بلکې د وت د
دوه‌چنده کولو دپاره بايد کار ته دوام ورکړي۔ را ياد کړئ چې دوه‌چنده
کوونکے ماشين د ننوت لومړۍ کرښه پاکوي او په بېل وت کښې دوه کرښې ليکي۔
راځئ يوه لارښوونه ور زياته کړو، څو ډاډمن شو چې د پټې سر د جمع ماشين
د وت لومړۍ کرښه لولي۔
\[
\begin{tikzpicture}[->,>=stealth',shorten >=1pt,auto,node distance=2.8cm,
                    semithick]
  \tikzstyle{every state}=[fill=none,draw=black,text=black]

  \node[initial,state] (A)              {$q_0$};
  \node[state]         (B) [right of=A] {$q_1$};
  \node[state]         (C) [right of=B] {$q_2$};
  \node[state]         (D) [below left of=C] {$q_3$};
  \node[state]         (E) [below left of=D] {$q_4$};

  \path (A) edge node {\TMtrans{\TMblank}{\TMstroke}{\TMstay}} (B)
            edge [loop above] node {\TMtrans{\TMstroke}{\TMstroke}{\TMright}} (B)
        (B) edge [loop above] node {\TMtrans{\TMstroke}{\TMstroke}{\TMright}} (B)
            edge node {\TMtrans{\TMblank}{\TMblank}{\TMleft}} (C)
        (C) edge node {\TMtrans{\TMstroke}{\TMblank}{\TMleft}} (D)
        (D) edge [loop left] node {\TMtrans{\TMstroke}{\TMstroke}{\TMleft}} (D)
            edge node[left, xshift=-2mm] {\TMtrans{\TMendtape}{\TMendtape}{\TMright}} (E);
\end{tikzpicture}
\]
اوس د ننوت دوه‌چنده کول اسانه دي---يوازې دوه‌چنده کوونکے ماشين له
حالت~$q_4$ سره نښلوو۔ د دې دپاره د دوه‌چنده کوونکي حالتونه داسې بيا
نومول پکار دي چې د~$q_0$ پر ځاے له~$q_4$ څخه پيل شي؛ په دې توګه دوه
لومړني حالتونه نۀ جوړېږي۔ وروستنی ډياګرام بايد د
\olref{fig:combined} په شان وي۔
\begin{figure}
\[
\begin{tikzpicture}[->,>=stealth',shorten >=1pt,auto,node distance=2.8cm,
                    semithick]
  \tikzstyle{every state}=[fill=none,draw=black,text=black]
  \node[initial,state] (A)              {$q_0$};
  \node[state]         (B) [right of=A] {$q_1$};
  \node[state]         (C) [right of=B] {$q_2$};
  \node[state]         (D) [below left of=C] {$q_3$};
  \node[state]         (E) [below left of=D] {$q_4$};
  \node[state]         (2) [right of=E] {$q_5$};
  \node[state]         (3) [right of=2] {$q_6$};
  \node[state]         (4) [below of=3] {$q_7$};
  \node[state]         (5) [left of=4]  {$q_8$};
  \node[state]         (6) [left of=5]  {$q_9$};

  \path (A) edge node {\TMtrans{\TMblank}{\TMstroke}{\TMstay}} (B)
            edge [loop above] node {\TMtrans{\TMstroke}{\TMstroke}{\TMright}} (B)
        (B) edge [loop above] node {\TMtrans{\TMstroke}{\TMstroke}{\TMright}} (B)
            edge node {\TMtrans{\TMblank}{\TMblank}{\TMleft}} (C)
        (C) edge node {\TMtrans{\TMstroke}{\TMblank}{\TMleft}} (D)
        (D) edge [loop left] node {\TMtrans{\TMstroke}{\TMstroke}{\TMleft}} (D)
            edge node[left, xshift=-2mm] {\TMtrans{\TMendtape}{\TMendtape}{\TMright}} (E)
    (E) edge node {\TMtrans{\TMstroke}{\TMblank}{\TMright}} (2)
    (2) edge [loop above] node {\TMtrans{\TMstroke}{\TMstroke}{\TMright}} (2)
      edge node {\TMtrans{\TMblank}{\TMblank}{\TMright}} (3)
    (3) edge [loop above] node {\TMtrans{\TMstroke}{\TMstroke}{\TMright}} (3)
        edge node {\TMtrans{\TMblank}{\TMstroke}{\TMright}} (4)
    (4) edge [loop below] node {\TMtrans{\TMblank}{\TMstroke}{\TMleft}} (4)
        edge node {\TMtrans{\TMstroke}{\TMstroke}{\TMleft}} (5)
    (5) edge [loop below]  node {\TMtrans{\TMstroke}{\TMstroke}{\TMleft}} (5)
        edge              node {\TMtrans{\TMblank}{\TMblank}{\TMleft}} (6)
    (6) edge [loop below] node {\TMtrans{\TMstroke}{\TMstroke}{\TMleft}} (6)
        edge              node {\TMtrans{\TMblank}{\TMblank}{\TMright}} (E);
\end{tikzpicture}
\]
\caption{د جمع کوونکي او دوه‌چنده کوونکي ماشينونو يوځاے کول}
\ollabel{fig:combined}
\end{figure}
\end{ex}

\begin{prop}
    که~$M$ او~$M'$ منضبط وي او په ترتيب سره تابعې $f\colon
    \Nat^k \to \Nat$ او $f'\colon \Nat \to \Nat$ محاسبه کوي، نو
    $M \frown M'$ منضبط دے او~$\comp{f}{f'}$ محاسبه کوي۔
\end{prop}

\begin{proof}
    څرنګه چې~$M$ منضبط دے، کله چې د وت
    ~$f(n_1,\dots,n_k) = m$ سره درېږي، سر مربع~$1$ لولي۔ که اوس د
    ~$M'$ لومړني حالت ته ننوځو، ماشين به د وت~$f'(m)$ سره ودرېږي او
    بيا به هم مربع~$1$ لولي۔ د~\olref[dis]{defn:disciplined} نور
    شرطونه هم پوره کېږي۔
\end{proof}

\begin{prob}
    د~\cref{tur:mac:dis:prob:disc-succ} ماشين~$M$ واخلئ،
    $M \frown M$ جوړ کړئ، او په دې توګه داسې منضبط ټيورينګ ماشين
    ورکړئ چې $f(x) = x+2$ محاسبه کوي۔
\end{prob}
\end{document}
